\section{v4.8 Research Integration}
\label{v480:integration}

This is a proposed integration layer for the research program named in
the repository, not a new release theorem. The complete v4.7 baseline
is retained below it. The layer separates what can be defined now from
what still requires a physical or mathematical construction.

\subsection{Typed Scale Record}
\label{v480:scale-record}

For a declared readout $R$, introduce the candidate scale record
\[
 \ell_{\rm eff}(R)=\ell_P\lambda_R,\qquad
 \lambda_R\geq0,
 \qquad \varepsilon_R^{\rm len}
       =\frac{\ell_P\lambda_R}{L_{\rm ref}},\quad L_{\rm ref}>0.
\]
Here $\ell_P$ and $L_{\rm ref}$ have units of length, and the new
$\lambda_R$ is dimensionless. The positive scale magnitude must be
kept separate from an orientation/signature marker $\sigma_R$.
In the inherited tachyon-facing passages a symbol such as
$\lambda_{\mathsf T}$ denotes a native scale magnitude, not
automatically this dimensionless coefficient. Identifying them
requires an explicit unit conversion and comparison map.

A signed interval class can therefore be recorded separately from
the length magnitude. A negative signature or orientation code
does not give a negative material length. Nor does a dimensionless
fit-score threshold such as $1/256$ equal
$\varepsilon_R^{\rm len}$ without a specified calibration between
the residual's norm and a length observable.

\subsection{Metric Ansatz and Null Limit}
\label{v480:metric-candidate}

If $g_{\rm ref}$ is a declared Lorentzian metric and $\lambda_R$
is smooth and strictly positive on a region, one possible additional
ansatz is
\[
 g_R=\lambda_R^2 g_{\rm ref}.
\]
Pointwise positive multiplication preserves signature and
nondegeneracy, and at a fixed point scales non-null interval
magnitudes by $\lambda_R$. This elementary property does not show
that the physical readout has this isotropic conformal form.
For a varying factor, connection and curvature terms also involve
derivatives of the factor; a length convention alone does not
determine their dynamics.

At $\lambda_R=0$ the proposed tensor is degenerate and is not a
Lorentzian metric. A null photon trajectory instead satisfies
$g(v,v)=0$ for a nonzero tangent in a nondegenerate metric; it does
not require the whole metric tensor to vanish. A photon-clock
boundary therefore needs its own null/readout construction, not
the substitution $\lambda_\gamma=0$ into a material metric ansatz.

The different null-preserving coframe construction in
Proposition~\ref{v480:hc-null-coframe} gives an explicit local geometric
example, not a derivation of this conformal ansatz. Its isolated stress
match is off shell. In particular, directly reusing its inverse coframe
in the same real-bilinear ordered master equation meets the exclusion
of Proposition~\ref{v480:hc-same-frame-obstruction}; microscopic and
effective frame data cannot be identified without this check.

\subsection{From MM/HC to Dynamics}
\label{v480:dynamics-handoff}

Before a metric or field equation can be claimed, specify the
state space and operator domains, the readout map and its units,
the represented MM and compensator, and the actual residual
functions. Proposition~\ref{v480:conditional-mm} states one
sufficient spectral interface; it does not construct the native
operator or its couplings.

The next dynamical package must also provide an evolution or action,
its boundary/initial conditions, the transformation of the readout
under admissible changes, and the required conservation or exchange
identities. Merely putting coefficient tags into geometric,
kinetic and mixed slots does not derive such a package.
The constant offset and nonlinear remainder in the C5 dictionary
must be retained when testing a finite contraction bound.

\subsection{Integration and Acceptance Gates}

\begin{enumerate}[leftmargin=*]
\item \textbf{Source and scope gate.} Keep the full-source map,
shifted normal form, backtranslation and terminal/global boundary.
Do not replace the later baseline with the older modular pilot.
\item \textbf{Scale gate.} Supply $R$, $L_{\rm ref}$, units,
positive magnitude, orientation/signature data and an evaluated
calibration before comparing length and detector residuals.
\item \textbf{Operator and residual gate.} Supply domains,
the spectral premises actually used and the native three-block
residual evaluation with its fixed norms. Synthetic regression
examples remain examples rather than native witnesses.
\item \textbf{Dynamical gate.} Derive the claimed equation from
the declared model and prove the needed well-posedness,
conservation, stability or approximation statements at the scope
being claimed. A fluid or gravitational interpretation is still
an interpretation until this map is constructed.
The declared real action, allowed variations and target reactions must
agree with the currents. The geometry/probe interface must also pass
the current and pointwise trace tests in
Section~\ref{v480:hc-dynamics-geometry}; an isolated Einstein/stress
match does not solve the missing probe and normal field equations.
Section~\ref{v490:cartan-source} partially advances this gate in an
explicitly added Cartan comparison. Its conditional subsidiary theorem
requires the full probe and constitutive equations on an open set;
its finite tangential-jet source law does not establish that hypothesis.
\item \textbf{Release gate.} Audit the corrected mathematical
statements and every dependent use, reconcile historical claims
with the current ledger, verify complete source inclusion and
layout, then make a deliberate version/DOI decision. No public
upload or exact all-readout claim follows from the present pass.
\end{enumerate}
